\documentclass[11pt]{article}

\usepackage[utf8]{inputenc}
\usepackage[T1]{fontenc}
\usepackage{geometry}
\usepackage{amsmath}
\usepackage{hyperref}
\usepackage{booktabs}

\geometry{margin=1in}

\title{Triangulation nach Schult:\\
A Structured Multi-Agent Protocol for Divergence Analysis and Human-Centered Synthesis}
\author{Denis Schult}
\date{April 2026}

\begin{document}

\maketitle

\begin{abstract}
The \textit{Triangulation nach Schult} is a structured, model-agnostic protocol for divergence analysis using multiple Large Language Models (LLMs). 
Instead of optimizing for consensus, the method treats disagreement, contradiction, and model refusal as epistemically informative signals.

The protocol defines a six-phase process (P1--P6) for hypothesis stress-testing, bias detection, and human-centered synthesis.
It is designed as a heuristic framework for decision-making under uncertainty, didactic applications, and exploratory system analysis.

The method does not aim to produce objective truth. Instead, it exposes weaknesses in assumptions through structured multi-perspective comparison.
An exploratory case study demonstrates its application using real model outputs.
\end{abstract}

\section{Introduction}

Large Language Models (LLMs) often produce convergent outputs due to shared training data and alignment constraints.
Such convergence may obscure systematic errors.

Traditional multi-model approaches emphasize agreement.
However, agreement is not a reliable proxy for correctness.

\textit{Triangulation nach Schult} reframes this:

\begin{quote}
Divergence is treated as an informative signal rather than noise.
\end{quote}

The method structures multiple LLM outputs to support human reasoning under uncertainty.

\section{Related Work}

The method relates to:

\begin{itemize}
\item Methodological triangulation (Denzin, Patton)
\item Delphi method (Dalkey \& Helmer)
\item Ensemble methods in machine learning
\item Adversarial testing and red teaming
\item Multi-agent LLM reasoning frameworks
\end{itemize}

Unlike these approaches, it does not optimize for consensus but analyzes structured divergence.

\section{Method Overview}

\subsection{Core Principle}

The method analyzes the structure of disagreement rather than seeking agreement.

\subsection{Functional Roles (Example)}

\begin{tabular}{ll}
\toprule
Role & Function \\
\midrule
Signal & Detects patterns and anomalies \\
Structure & Builds relationships and scenarios \\
Critique & Identifies logical weaknesses \\
Boundary & Marks uncertainty and limits \\
Falsification & Tests hypotheses \\
Synthesis & Integrates outputs \\
Human Operator & Final judgment \\
\bottomrule
\end{tabular}

\section{Protocol (P1--P6)}

\textbf{P1:} Define hypotheses and falsification thresholds.\\
\textbf{P2:} Select models and assign roles.\\
\textbf{P3:} Parallel prompting with identical input.\\
\textbf{P4:} Map divergence across outputs.\\
\textbf{P5:} Human synthesis with justification.\\
\textbf{P6:} Meta-reflection on bias and omissions.

\section{Falsification Conditions}

\begin{itemize}
\item Persistent divergence without coherent interpretation
\item Insufficient model coverage
\item Unjustified synthesis
\item External contradiction
\end{itemize}

\section{Case Study (Exploratory)}

\subsection{Setup}

On April 4, 2026, multiple LLMs were prompted with an identical query:

\begin{quote}
Analyze the European energy situation (April 2026) focusing on TTF gas prices and industrial response.
Evaluate predefined hypotheses and provide structured reasoning.
\end{quote}

Models included:
\begin{itemize}
\item Grok (signal-focused)
\item Claude (analytical reasoning)
\item Copilot (operational perspective)
\end{itemize}

\subsection{Observations}

\begin{itemize}
\item Grok emphasized high volatility and potential cascading effects.
\item Claude identified structural industrial adaptation and long-term constraints.
\item Copilot declined full evaluation due to insufficient hypothesis specification.
\end{itemize}

\subsection{Interpretation}

The divergence revealed three distinct layers:

\begin{itemize}
\item Empirical signals (market volatility)
\item Structural interpretation (industrial adaptation)
\item Methodological constraint (insufficient hypothesis definition)
\end{itemize}

The refusal of one model functioned as a meta-signal, indicating incomplete operationalization of hypotheses.

\subsection{Outcome}

The synthesized result was a cautious interpretation:
no immediate investment decision recommended due to unresolved uncertainty and incomplete hypothesis specification.

This case illustrates:

\begin{itemize}
\item Divergence as a diagnostic signal
\item Model refusal as methodological feedback
\item Structured synthesis under uncertainty
\end{itemize}

This does not constitute formal validation but demonstrates feasibility.

\section{Didactic Implications}

The method supports:

\begin{itemize}
\item perspective comparison
\item bias awareness
\item reflective reasoning
\end{itemize}

It is suitable for vocational and higher education contexts.

\section{Limitations}

\begin{itemize}
\item No validation of external reality
\item Dependence on model quality
\item Sensitivity to hypothesis formulation
\item Human interpretation required
\end{itemize}

\section{Conclusion}

The method provides a structured way to analyze divergence in LLM outputs.
It does not generate truth but exposes uncertainty and supports justified decision-making.

\section{Availability}

Released as Version 1.1-Open.

\section{Citation}

Schult, D. (2026). 
Triangulation nach Schult: A Structured Multi-Agent Protocol for Divergence Analysis and Human-Centered Synthesis.

\end{document}
